\section*{Chapter \ref{ch:g9:inside-the-atom} --- Inside the Atom}

\begin{solution}{exo:g9:inside-the-atom:1}
In the \omterm{def:g9:inside-the-atom:nucleus}{nucleus}: \omterm{def:g9:inside-the-atom:proton}{protons} (positive) and \omterm{def:g9:inside-the-atom:proton}{neutrons} (no charge). Around it:
\omterm{def:g9:inside-the-atom:nucleus}{electrons} (negative).
\end{solution}

\begin{solution}{exo:g9:inside-the-atom:2}
\ce{^{16}_{8}O}: 8 \omterm{def:g9:inside-the-atom:proton}{protons}, 8 \omterm{def:g9:inside-the-atom:proton}{neutrons}, 8 \omterm{def:g9:inside-the-atom:nucleus}{electrons}.
\ce{^{27}_{13}Al}: 13, 14, 13. \ce{^{40}_{20}Ca}: 20, 20, 20.
\end{solution}

\begin{solution}{exo:g9:inside-the-atom:3}
It has as many \omterm{def:g9:inside-the-atom:nucleus}{electrons} (charge $-e$ each) as \omterm{def:g9:inside-the-atom:proton}{protons} (charge $+e$
each): the charges cancel.
\end{solution}

\begin{solution}{exo:g9:inside-the-atom:4}
$A = 26 + 30 = 56$: \ce{^{56}_{26}Fe}.
\end{solution}

\begin{solution}{exo:g9:inside-the-atom:5}
Yes: both have $Z = 17$, so both are chlorine. They have different \omterm{def:g9:inside-the-atom:atomic-number}{mass
numbers} (18 and 20 \omterm{def:g9:inside-the-atom:proton}{neutrons}): they are \omterm{def:g9:inside-the-atom:isotope}{isotopes}.
\end{solution}

\begin{solution}{exo:g9:inside-the-atom:6}
No: they have different \omterm{def:g9:inside-the-atom:atomic-number}{atomic numbers} (6 and 7), so they are different
elements, carbon and nitrogen. \omterm{def:g9:inside-the-atom:isotope}{Isotopes} have the same $Z$.
\end{solution}

\begin{solution}{exo:g9:inside-the-atom:7}
$56 \times 1.67 \times 10^{-27} \approx \qty{9.35e-26}{kg}$.
\end{solution}

\begin{solution}{exo:g9:inside-the-atom:8}
An \omterm{def:g7:atoms-and-molecules:atom}{atom} is mostly empty space: the \omterm{def:g9:inside-the-atom:nucleus}{nucleus} is tiny, so most particles
meet no \omterm{def:g9:inside-the-atom:nucleus}{nucleus} on their way and pass straight through.
\end{solution}

\begin{solution}{exo:g9:inside-the-atom:9}
$10^{-10} / 10^{-15} = 10^{5}$: a hundred thousand times smaller.
\end{solution}

\begin{solution}{exo:g9:inside-the-atom:10}
Chemistry depends on the \omterm{def:g9:inside-the-atom:nucleus}{electrons}, and \omterm{def:g9:inside-the-atom:isotope}{isotopes} of an element have the
same number of \omterm{def:g9:inside-the-atom:nucleus}{electrons} (the same $Z$).
\end{solution}

\begin{solution}{exo:g9:inside-the-atom:11}
$26 \times 9.11 \times 10^{-31} = \qty{2.37e-29}{kg}$. Fraction:
$2.37 \times 10^{-29} / 9.35 \times 10^{-26} \approx 2.5 \times
10^{-4}$, about \qty{0.025}{\percent} of the mass.
\end{solution}

\begin{solution}{exo:g9:inside-the-atom:12}
6915 copper-63 \omterm{def:g7:atoms-and-molecules:atom}{atoms} and 3085 copper-65 \omterm{def:g7:atoms-and-molecules:atom}{atoms}. Average:
$(6915 \times 63 + 3085 \times 65)/10\,000 = 63.6$ \omterm{def:g9:inside-the-atom:proton}{nucleons}.
\end{solution}

\begin{solution}{pb:g9:inside-the-atom:1}
\textbf{1.} \ce{^{35}_{17}Cl} and \ce{^{37}_{17}Cl}.

\textbf{2.} Chlorine-35: 17 \omterm{def:g9:inside-the-atom:proton}{protons}, 18 \omterm{def:g9:inside-the-atom:proton}{neutrons}, 17 \omterm{def:g9:inside-the-atom:nucleus}{electrons}.
Chlorine-37: 17 \omterm{def:g9:inside-the-atom:proton}{protons}, 20 \omterm{def:g9:inside-the-atom:proton}{neutrons}, 17 \omterm{def:g9:inside-the-atom:nucleus}{electrons}.

\textbf{3.} Both have 17 \omterm{def:g9:inside-the-atom:proton}{protons}, so both are the element chlorine; they
differ only by their number of \omterm{def:g9:inside-the-atom:proton}{neutrons}, so they are \omterm{def:g9:inside-the-atom:isotope}{isotopes}.

\textbf{4.} Yes: chemistry depends on the \omterm{def:g9:inside-the-atom:nucleus}{electrons}, and both have 17.

\textbf{5.} $35 \times 1.67 \times 10^{-27} = \qty{5.845e-26}{kg}$.

\textbf{6.} $17 \times 9.11 \times 10^{-31} = \qty{1.55e-29}{kg}$:
about 4000 times less than the \omterm{def:g9:inside-the-atom:nucleus}{nucleus}, negligible.

\textbf{7.} $37 \times 1.67 \times 10^{-27} = \qty{6.179e-26}{kg}$.

\textbf{8.} $\qty{1}{g} = \qty{e-3}{kg}$; $10^{-3} / 5.845 \times
10^{-26} \approx 1.71 \times 10^{22}$ \omterm{def:g7:atoms-and-molecules:atom}{atoms}.

\textbf{9.} \qty{75.8}{\percent} and \qty{24.2}{\percent}.

\textbf{10.} $758 \times 5.845 \times 10^{-26} = \qty{4.430e-23}{kg}$;
$242 \times 6.179 \times 10^{-26} = \qty{1.495e-23}{kg}$.

\textbf{11.} $4.430 \times 10^{-23} + 1.495 \times 10^{-23} =
\qty{5.925e-23}{kg}$.

\textbf{12.} $5.925 \times 10^{-23} / 1000 \approx
\qty{5.93e-26}{kg}$: the average mass of a chlorine \omterm{def:g7:atoms-and-molecules:atom}{atom}.
\end{solution}
